\section{Guile}

\textbf{TODO:
\begin{itemize}
  \item{Historical record of variadic functions}
  \item{Talk about lambda-case in background (maybe, this is an intermediate representation which I considered out-of-scope for Python)}
  \item{Either include bytecode in the destructuring section or explain why you don't}
  \item{Explain distinction between Guile, Scheme, and C}
  \item{Discuss rNrs in background \& particularly r2rs}
  \item{Remove most of the words from the background section...}
  \item{Fix citation to manual to properly cite info manual instead of repo url}
  \item{Figure out citation for devel archive}
  \item{Double-check bytecode samples are still correct after aux refactoring}
  \item{Normalize casing of "c-level"}
  \item{Normalize bytecode explanations (refer to specific instruction numbers, be consistent about whether the explanation or the snippet comes first)}
  \item{Normalize manner of discussing cross-feature implications. Inline, footnotes, or dedicated section?}
  \item{Add explanation of subject of study to introduction}
\end{itemize}
}

Definition of define*
\parencite[][modules/system/ice-9/psyntax.scm lines 3377-3383]{guile-source}.

Guile provides the following features:

\begin{itemize}
  \item{Destructuring}
  \item{Named Value}
  \item{Default Values}
  \item{Positional Value}
  \item{Typed Values}
  \item{Variadic Functions}
\end{itemize}

\subsection{Background}
Guile is an implementation of the scheme programming language, which itself is a dialect
of the lisp programming language.
Lisp, which stands for "LISt Processing" is a family of languages which are grouped
together because of their relationship to a programming language originally developed by
a small group of people with support from various institutions including the
Massachussettes Institute of Technology and the United States military
\parencite[][page 195]{original-lisp-paper-pt-1}.
The most immediately noticable differences from common languages like C++ and Python is
that they use prefix notation, put the open parenthesis before the function name, and
allow a wider variety of characters to appear in variable names.
For example, this is how one would call a getter function in lisp, comparable to
\texttt{object\_variable.getter\_func\_name()}:

\begin{pygmented}[lang=lisp]
(getter-func-name object-variable)
\end{pygmented}

Note that the function name comes before the object name, the parentheses enclose both
names, and hyphens are used instead of underscores.

Prefix notation applies to primitive mathematical operations too.
For example, one can add 3 numbers in lisp with this expression:

\begin{pygmented}[lang=lisp]
(+ 0 1 2)
\end{pygmented}

It implements a compiler which produces bytecode, an interpreter which executes bytecode,
and an interpreter which directly executes program text
\parencite[][modules/language/scheme/spec.scm lines 43-45]{guile-source}.
This paper will focus on compiled bytecode as this keeps the analysis consistent with
other sections and is the typical way to execute Guile code\footnote{
  For example, running \texttt{guile script-name} will first compile the script then run
  the compiled file rather than running the script directly.
}.

Guile bytecode operates as a stack machine with 2 pointers into the stack: the frame
pointer and the stack pointer.
\parencite[][Section 9.3.2]{guile-manual}.
The frame pointer stores a location near\footnote{
  It is "near" the beginning rather than "at" the beginning because some metadata is
  stored before the frame pointer's location; this metadata will not be relevant to this
  paper.
}
the beginning of the frame, where each frame represents a single function call.
The stack pointer keeps track of the end of the stack, like a pointer to the end of an
array
\parencite[][Section 9.3.3]{guile-manual}.
When instructions take an index, they differ in whether they take in an index relative to
the stack pointer or the frame pointer
\parencite[][Section 9.3.5]{guile-manual}.
This paper will always reference indexes relative to the stack pointer for the sake of
consistency.

\subsubsection{The optargs module}
Many of the features described in this paper are provided through Guile's \texttt{optargs}
module.
Here I will explain the general syntax of the \texttt{optargs} module and its history.

\textbf{Syntax}

In plain scheme, functions are defined like this:

\begin{pygmented}[lang=lisp]
(define (add-values a b)
  (+ a b))
\end{pygmented}

The function's name is \texttt{add-value} and the arguments are \texttt{a} and \texttt{b}.
Callers can only call the function by providing positional values for every argument:

\begin{pygmented}[lang=lisp]
(add-values 2001 2002)
;; results in 3003
\end{pygmented}

The \texttt{optargs} module adds a new way of defining functions, with \texttt{define*}.
This can take the same inputs that a plain \texttt{define} takes, in which case the
semantics are identical:

\begin{pygmented}[lang=lisp]
(define* (add-values a b)
  (+ a b))
\end{pygmented}

Authors can change the semantics by adding keywords to the function signature.
Keywords modify all following arguments, until a different keyword is reached.
For example, this definition marks both \texttt{a} and \texttt{b} as "optional arguments",
which are described in the section on default values:

\begin{pygmented}[lang=lisp]
(define* (add-values #:optional a b)
  (+ a b))
\end{pygmented}

This one marks both as "keyword arguments", which are described in the section on named
values:

\begin{pygmented}[lang=lisp]
(define* (add-values #:key a b)
  (+ a b))
\end{pygmented}

This one marks \texttt{a} as an optional argument and \texttt{b} as a keyword argument.
Note that \texttt{b} is \textit{not} an optional argument; it is only a keyword argument
(though as you will see in the following sections, their semantics partially overlap):

\begin{pygmented}[lang=lisp]
(define* (add-values #:optional a #:key b)
  (+ a b))
\end{pygmented}

\textbf{History}

The \texttt{optargs} module was introduced in version 1.3.2, released in 1999
\parencite[][Section 9.1.4 "A Timeline of Selected Guile Releases"]{guile-manual}
\parencite[][NEWS lines 11451-11525]{guile-source}. 
At this time, Guile lacked a virtual machine and the module was implemented in pure
scheme.
This implementation operated by adding a prelude to the function definiton which searches
through the values provided by the caller to decide which values should be assigned to
which variables
\parencite[][ice-9/optargs.scm]{guile-1.3.2}.
The virtual machine was added to Guile in version 2.0.0, released in 2010
\parencite[][Section 9.1.4 "A Timeline of Selected Guile Releases"]{guile-manual}
\parencite[][NEWS lines 5068-5071]{guile-source}.
This also entailed a rewrite of the \texttt{optargs} module which centered the
implementation around the internal \texttt{<lambda-case>} strucure.
This structure contains all the information about the different kinds of arguments
accepted by the function and facilitates the use of special-purpose VM instructions for
processing arguments efficiently
\parencite[][2009-10 lines 6428-6453]{guile-devel-archive}.
While the VM has gone through changes since then, including a complete rewrite in version
2.2
\parencite[][Section 9.1.4 "A Timeline of Selected Guile Releases"]{guile-manual},
the implementation of the \texttt{optargs} module has remained stable.
The special-purpose VM instructions have been adapted to handle the details of VM
operation correctly and decrease latency, but the core logic they use to process arguments
has proven to be robust.

7c54029740a147a623c1f0564708d5471addf232 (Oct 03 2013) through
553294d958c953f57658bad45affc15b55fcc471 (Nov 28 2013) - VM Rewrite

df1cd5e59b38a7ecdb2184d7f0fed3f803969541 - (Oct 23 2009) Moves lambda* to psyntax
4d3406a84750589ffa70669bb923343082efaeb5 - (Oct 23 2009) optargs rewrite
a16d4e82e949954805bf2cd42cfbb519fcf4012d - (Apr 06 2013) Adds bind-kwargs
1c33be992e8120abd20add8021e4d91d226f5b6a - (Nov 08 2013) Removes the "old VM"
8e5755e7719b2bab69c509e4bf0ab2e8bcdc8a10 - (Jun 24 2018) Adds bind\_kwargs intrinsic.

\subsection{Simple Function and Call}

\begin{pygmented}[lang=lisp]
(define (add-values a b)
  (+ a b))

(add-values 1000 1001)
\end{pygmented}

The bytecode
\footnote{
	The partial-evaluation optimization was turned off for all of the samples given here.
	This is because partial evaluation reduces the entire program to a single number.
}
generated for the call to \texttt{add-values} performs the function call by loading
the function, loading the arguments, and dispatching to the function.
Instructions 0-2 load the \texttt{add-values} procedure into the beginning of the stack
and instruction 3 moves it into index 5.
Instructions 4-5 place the constant values into the 4th and 3rd indices.
Instruction 6 can be safely ignored; it handles events like signals from the operating
system.
Finally, instruction 7 calls the function.
This involves moving the frame pointer to index 5 (where the procedure is stored) and the
stack pointer to index 2 (the new end of the stack, which is smaller for the callee than
for the caller).

\begin{pygmented}[lang=lisp]
0 (call-scm<-thread 0 62)         ;; current-module
1 (static-ref 1 16322)            ;; add-values
2 (call-scm<-scm-scm 0 0 1 111)   ;; lookup-bound
3 (scm-ref/immediate 5 0 1)       
4 (make-immediate 4 4002)         ;; 1000
5 (make-immediate 3 4006)         ;; 1001
6 (handle-interrupts)             
7 (call 4 3)                      
\end{pygmented}

The bytecode for the function definition is similar but simpler.
It calls the built-in function \texttt{add} (which corresponds to the \texttt{+}
function).
Because \texttt{add} is a built-in function it does not need to be loaded or placed on
the stack and its arguments can come from any stack location.

\begin{pygmented}[lang=lisp]
0 (call-scm<-scm-scm 7 6 5 0)     ;; add
1 (reset-frame 1)                 ;; 1 slot
2 (handle-interrupts)             
3 (return-values)                 
\end{pygmented}

\subsection{Caller Destructuring}
Provided through the \texttt{apply} function. This allows a caller to provide
a list which is destructured into a set of positional arguments.

\subsubsection{Syntax \& Semantics}
Destructuring in guile has a simple interface, as \texttt{apply} is simply a function that
takes in a function as the first value, then any number of arbitrary values, and requires
a possibly empty list as the final value.
For example, all of these are equivalent:

\begin{pygmented}[lang=lisp]
(add-values 3 4)
(apply add-values 3 '(4))
(apply add-values '(3 4))
\end{pygmented}

However, this would be an error because the final value is not a list:

\begin{pygmented}[lang=lisp]
(apply add-values 3 4)
\end{pygmented}

\subsubsection{Implementation}
\texttt{Apply} operates by first collecting all of the positional arguments into a single
list, with the final value as the tail of the list, and sending them to
\texttt{scm\_call\_n}
\parencite[][libguile/eval.c lines 585-622, lines 715-729]{guile-source}
\footnote{
  A typical guile programmer (such as myself) who is familiar with the dotted list syntax
  or the \texttt{\#:rest} argument (both described in the section on variadic functions)
  might be confused by the c-level API of the apply function.
  It takes 2 required arguments and a rest argument.
  The apply function takes any number of arguments and requires that the final argument is
  a list.
  However, the body of the function simply calls \texttt{cons*} on the second required
  argument and the rest argument; one might expect that this would cause the user-provided
  list to be preserved instead of flattened.
  However, when a c-level function is registered with a rest argument in guile it receives
  an improper list of arguments, rather than the proper list provided when
  \texttt{\#:rest} is given to \texttt{define*}.
  See the file \texttt{examples/guile/c-gsubr-rest.scm} (and its corresponding c file) in
  the paper's repository for a demonstration.
  The reason for this discrepency is not clear but the API of \texttt{scm\_apply} makes
  sense now.
}.
This function is responsible for setting up the VM stack correctly
\parencite[][libguile/vm.c lines 1542-1621]{guile-source},
a task which would normally be handled by the \texttt{call} instruction
\parencite[][libguile/vm-engine.c lines 428-461]{guile-source}.

\subsubsection{Historical Record}
The notion of list destructuring with \texttt{apply} was included in the original paper
defining the lisp programming language
\parencite[][pages 189-190]{original-lisp-paper-pt-1}.
This original definition took exactly 2 parameters, a function and a list containing the
arguments.
The second revised report on scheme extended this by allowing the caller to pass any
number of values between the function and the list
\parencite[][page 57]{r2rs}.

Guile's original implementation did not have a virtual machine; it contained only an
interpreter that worked with a literal representation of the program text
\parencite[][Section 9.3.1 "Why a VM?"]{guile-manual}.
The earliest commit in the current guile repository simply called the function after
manually unpacking the arguments
\parencite[][libguile/eval.c lines 1869-1985]{guile-first-commit}.
With the release of Guile 2.0, which added the virtual machine
\parencite[][NEWS lines 5068-5071]{guile-source},
the \texttt{apply} function delegated to \texttt{scm\_call\_with\_vm}, a dedicated c-level
helper function for calling scheme-level functions
\parencite[][libguile/eval.c lines 782-802]{guile-2.0}.
At this point, the VM was implemented as a C-level function which executed a single
C-level function \textbf{TODO: I'm not sure if this sentence is worded confusingly enough
yet...}, so the end result was simply to call this function with the appropriate
arguments
\parencite[][libguile/vm.c lines 560-566, libguile/vm-engine.c lines 36-296]{guile-2.0}.
This was replaced with the \texttt{scm\_call\_n} function as part of a refactoring which
removed explicit references to the VM from many locations in C code as well as eliminating
VM visibility from scheme code
\parencite[][commit 165d9bf3a3bf34b53ed916743c6414f8030320c3 and earlier commits found
with \texttt{git log libguile/eval.c}]{guile-source}.
The reason for this change is not explicitly stated in the commit messages but is likely
related to performance improvements attributed to a rewrite of the virtual machine in the
Guile 2.2 release notes
\parencite[][NEWS lines 1995-2004, 2124-2205]{guile-source}.

\subsection{Author Destructuring}
This feature is not provided by Guile.

\subsection{Implicit Value}

This feature is not provided by Guile.

\subsection{Named Value}
Provided through the \texttt{\#:key} arguments to \texttt{optargs} definitions.

\subsubsection{Syntax \& Semantics}
The caller must use the name specified by the author, which will also be the name used by
the local variable
\footnote{
	The caller uses a keyword to name the value while the author uses a symbol to name the
	variable, as normal, but they both use the same name.
}.
When an author mandates a named value they also implicitly create a default value.
Details about default values are described in the appropriate section; for the purpose of
named values all we need to know is that if a caller omits a value then the local variable
will be bound to \texttt{\#:false}.
For example, an author of \texttt{add-values} could allow named values like this:

\begin{pygmented}[lang=lisp]
(define* (add-values #:key a b)
  (+ (or a 0) (or b 0)))
\end{pygmented}

This function could be called in any of the following ways, with the result shown in the
proceeding comment:

\begin{pygmented}[lang=lisp]
(add-values)             ; 0
(add-values #:a 3)       ; 3
(add-values #:b 3)       ; 3
(add-values #:a 3 #:b 3) ; 6
\end{pygmented}

However, these would not be legal:

\begin{pygmented}[lang=lisp]
(add-values 3)
(add-values 2 #:b 2)
\end{pygmented}

When an author specifies that a value \textit{can} be named this also means that it
\textit{must} be named, if it is provided at all.

Additionally, an author may specify that a caller can supply arbitrary named values which
are not specified in the function signature (and not bound to a local variable, unless
the function is also variadic).
This is done by adding \texttt{\#:allow-other-keys}
\footnote{
	Ironically, while the optargs module does not support implicit values for user-defined
	code, \texttt{\#:allow-other-keys} is itself an implicit value.
}
to the function signature.
For example, a definition of \texttt{add-values} with this feature would allow all of
the following calls to be legal, and they would all produce the value 6.

\begin{pygmented}[lang=lisp]
(define* (add-values #:key (a 0) (b 0)
                     #:allow-other-keys)
  (+ (or a 0) (or b 0)))

(add-values #:a 3 #:b 3)
(add-values #:a 3 #:b 3 #:c 34)
(add-values #:c 34 #:a 3 #:b 3)
(add-values #:a 3 #:b 3 #:multiplier 88)
\end{pygmented}

\subsubsection{Implementation}

When Guile compiles a procedure it stores the list of valid argument names next to the
\texttt{SCM} object representing the code
\parencite[][libguile/vm-engine.c lines 734, 736, 743-745;
             libbugile/vm.c       lines 1003, 1008-1009]{guile-source}.
\textbf{TODO: 6.7.3 is the scheme-level API for accessing procedure metadata}
The body of the \texttt{add} function compiles to the following bytecode:

\begin{pygmented}[lang=lisp]
0  (bind-kwargs
     1 0 1 3 16367)
1  (alloc-frame 3)       ;; 3 slots
2  (immediate-tag=?
     1 4095 2308)        ;; undefined?
3  (jne 2)               ;; -> L1
4  (make-immediate 1 14) ;; 3
  
5  (immediate-tag=?
     0 4095 2308)        ;; undefined?
6  (jne 2)               ;; -> L2
7  (make-immediate 0 4)  ;; #f

8  (call-scm<-scm-scm
      2 1 0 0)           ;; add
9  (reset-frame 1)       ;; 1 slot
10 (handle-interrupts)
11 (return-values)
\end{pygmented}

The first instruction, \texttt{bind-kwargs}, performs the bulk of the work required to
process named values at runtime.
First, it initializes all relevant local variables (eg, ones that have default values or
are named) to the undefined value.
Next, it walks through the list of non-positional arguments with the assumption that every
even-indexed item is a keyword naming an argument
\footnote{
  Technically, it assumes that alternating value are keywords, except that it will
  silently ignore a non-keyword if \texttt{\#:rest} is included in the procedure
  definition.
  However, I am currently focused on named values in isolation of other features.
  \texttt{\#:rest} is discussed in the section on variadic functions.
}
It is an error if a value which should be a keyword is not or if a keyword does not name
an argument, unless \texttt{\#:allow-other-keys} was specified.
\parencite[][libguile/vm.c lines 977-1039]{guile-source}.
The function's bytecode includes instructions to fill variables with their default values
if they are not given an explicit value; this is discussed further in the section on
default values (instructions 2-7).

The bytecode for the call is largely similar to the simple call, except that it also
pushes the keywords onto the stack:

\begin{pygmented}[lang=lisp]
(call-scm<-thread 7 62)         ;; current-module
(static-ref 6 16329)            ;; sum
(call-scm<-scm-scm 7 7 6 41)    ;; define!
(make-non-immediate 6 16327)    ;; #<procedure sum (#:key a b)>
(scm-set!/immediate 7 1 6)      
(static-ref 3 16338)            ;; #:a
(make-immediate 2 34)           ;; 8
(static-ref 1 16347)            ;; #:b
(make-immediate 0 38)           ;; 9
(make-immediate 4 4)            ;; #f
(handle-interrupts)
(call-label 3 5 8)              ;; sum at #x7f308dda2158
\end{pygmented}

\subsubsection{Historical Record}
Named values were introduced with the original implementation of the \texttt{optargs}
module \parencite[][NEWS lines 11451-11525]{guile-source}. 
In the original implementation, the \texttt{let-ok-template} helper function generated a
scheme-level \texttt{let} (or \texttt{let*}) form which initially binds each variable name
to either the default value supplied by the author or a fresh undefined variable
\parencite[][commit 7e01997e88c54216678271de36b1c2088377492d ice-9/optargs.scm lines 128-135]{guile-source}.
The \texttt{bindfilter} local in the \texttt{let-keyword-template} helper replaced these
initial values with the values supplied by the caller
\parencite[][commit 7e01997e88c54216678271de36b1c2088377492d ice-9/optargs.scm lines 152-168]{guile-source}.

The current implementation is largely the same, except that it takes advantage of
being able to directly manipulate the VM state from C code instead of using macros to
mutate code syntactically.
The maintainers decided to move the implementation from scheme to C in order to decrease
the latency associated with a function call that uses named values
\parencite[][2009-10 lines 10053-10056, 14008-14009, 14139-14141]{guile-devel-archive}.

One change, which may or may not be considered a bug fix, has to do with the way that
Guile processes named values when the function also contains a variadic variable.
Originally, Guile required that all named values precede all (other) variadic values.
This was changed so that variadic values may be intermingled with named values, so long
as variadic values are not keywords (unless \texttt{\#:allow-other-keys} is set, in which
case they may be keywords)
\parencite[][commit ff74e44ecba55f50b2c2c84bad2f13bed9489455]{guile-source}.

\subsection{Default Value}
Provided through both \texttt{\#:optional} and \texttt{\#:key}.
When considering default values, both of these are semantically equivalent.
The difference between them is whether values are passed by position or by name.
This section will exclusively use \texttt{\#:optional} because none of the functions used
in the examples are complicated enough to benefit from named values.

\subsubsection{Syntax \& Semantics}
The default value for an argument can either be specified or unspecified.
If it is unspecified it is \texttt{\#false}.
Arguments become associated with a default value by placing them after the
\texttt{\#:optional} keyword in a \texttt{define*} form.
This verison of \texttt{add-values} demonstrates use of default values without
specification.

\begin{pygmented}[lang=lisp]
(define* (add-values #:optional a b)
  (+ (or a 0) (or b 0)))
\end{pygmented}

This function adds the given values, and the body of the function uses the value
\texttt{0} if the caller omitted a vaule.

However, it is idiomatic to specify a default value in the function signature rather than
testing the truthiness of the value.
This is done by specifying the argument with a list instead of a symbol.
The first element is the symbol naming the local variable and the second element is the
default value.

\begin{pygmented}[lang=lisp]
(define* (add-values #:optional (a 0) (b 0))
  (+ a b))
\end{pygmented}

This version is semantically equivalent to the first.

\subsubsection{Implementation}
The bytecode for calling a function using default values does not differ from the bytecode
for the simple call.

The bytecode for the function definition differs significantly.
I will start with the bytecode for the function which includes default values in the
signature as this version is simpler:

\begin{pygmented}[lang=lisp]
0  (bind-optionals 3)              ;; 2 args
1  (alloc-frame 8)                 ;; 8 slots
2  (immediate-tag=? 6 4095 2308)   ;; undefined?
3  (jne 2)                         ;; -> L1
4  (make-immediate 6 2)            ;; 0
L1:
5  (immediate-tag=? 5 4095 2308)   ;; undefined?
6  (jne 2)                         ;; -> L2
7  (make-immediate 5 2)            ;; 0
L2:
8  (call-scm<-scm-scm 7 6 5 0)     ;; add
9  (reset-frame 1)                 ;; 1 slot
10 (handle-interrupts)
11 (return-values)
\end{pygmented}

The first instruction, \texttt{bind-optionals} checks if the caller omitted any argument
values.
If so, it fills in the associated variables with the undefined value
\parencite[][libguile/vm-engine.c lines 3213-3233]{guile-source}.
The next 2 sections (instructions 2-4 and 5-7) check whether or not the optional values
are undefined; if so, they are filled with the default value specified by the author.
Finally, instructions 8-11 call the \texttt{add} function and return it's value
identically to the simple call.

The version which does not specify default values in the signature is similar.
It also calls \texttt{bind-optionals} (instruction 0) then fills in a default value if
any of the arguments are undefined (instructions 2-7).
However, in this case the default value is \texttt{\#:false}.
The next two sections check the argument values a second time and replace the argument
value with \texttt{0} if it was false (instructions 8-14).
Finally, it calls the \texttt{add} function and returns it's value identically to the
simple call (instructions 15-18).

\begin{pygmented}[lang=lisp]
0  (bind-optionals 3)              ;; 2 argss
1  (alloc-frame 3)                 ;; 3 slots
2  (immediate-tag=? 1 4095 2308)   ;; undefined?
3  (jne 2)                         ;; -> L1
4  (make-immediate 1 4)            ;; #f
L1:
5  (immediate-tag=? 0 4095 2308)   ;; undefined?
6  (jne 2)                         ;; -> L2
7  (make-immediate 0 4)            ;; #f
L2:
8  (immediate-tag=? 1 3839 4)      ;; false?
9  (jne 2)                         ;; -> L3
10 (make-immediate 1 2)            ;; 0
11 L3:
12 (immediate-tag=? 0 3839 4)      ;; false?
13 (jne 2)                         ;; -> L4
14 (make-immediate 0 2)            ;; 0
L4:
15 (call-scm<-scm-scm 2 1 0 0)     ;; add
16 (reset-frame 1)                 ;; 1 slot
17 (handle-interrupts)             
18 (return-values)                 
\end{pygmented}

\subsubsection{Historical Record}
Default values were included in the original \texttt{optargs} implementation
\parencite[][NEWS lines 11451-11525]{guile-source}. 
The semantics at that time were nearly identical to the current semantics.
The one notable difference is that, in the case where the author does not define a
default value and the caller does not supply one, the variable was not bound to any value.
This was changed as part of merging a branch that was mostly concerned with how modules
are handled.
I was unable to find any record about the reason why this change was made
\footnote{
	Marius Vollmer, the developer who authored the commit, was very kind in responding to an
	email enquiring about this.
	Unfortunately it has been more than 2 decades since the change was made and it was not
	notable enough to warrant a place in permanent memory.
}
\parencite[][commit 296ff5e78b8322fe4bf00c5ec1497dc28da776b8]{guile-source}
\parencite[][2001-05 lines 11305-11313]{guile-devel-archive}.
This was included in the 1.8 release series first released in 2006
\parencite[][Section 9.1.4 "A Timeline of Selected Guile Releases"]{guile-manual}
\footnote{
	I infer that it was first added in 1.8 because commit
	296ff5e78b8322fe4bf00c5ec1497dc28da776b8, which implements the change, is between
	commit 0f24e75b73b9c0c745971de639a53748a395a1cb, which bumps the numbers in
	GUILE-VERSION from 1.7 to 1.9 (guile uses odd numbers for development versions and even
	numbers for release versions) and commit c299f186ff9693fc88859daef037e3d94cc7c0ff, which
	adds content to the NEWS file regarding changes new in 1.6
	However, I did not find any content in the NEWS file which discusses the change
	(including in more recent release notes).
}.

\subsection{Positional Value}
Provided as the default mechanism for mapping values to arguments.

\subsubsection{Syntax \& Semantics}
The syntax and semantics for both calls and definitions are identical to the simple call.

\subsubsection{Implementation}
Guile does not explicitly enforce a restriction that values can only be supplied
positionally.
This is because, technically, all values are always supplied positionally.
When a caller provides named values they are actually providing pairs of values: keywords
as names and arbitrary values as values.
For example, given this call it is impossible to say whether the caller is supplying
named values or positional values:

\begin{pygmented}[lang=lisp]
(make-dictionary #:first  "one"
                 #:second "two"
                 #:third  "three")
\end{pygmented}

This could be a valid call to a function with named values:

\begin{pygmented}[lang=lisp]
(define* (make-dictionary #:key first second third)
  (let ((result (make-hash-table)))
    (hash-set! result #:first  first)
    (hash-set! result #:second second)
    (hash-set! result #:third  third)
    result))
\end{pygmented}

Or it could be a valid call to a function which takes 6 positional values:

\begin{pygmented}[lang=lisp]
(define* (make-dictionary key0 val0
                          key1 val1
                          key2 val2)
  (let ((result (make-hash-table)))
    (hash-set! result key0 val0)
    (hash-set! result key1 val1)
    (hash-set! result key2 val2)
    result))
\end{pygmented}

From the perspective of this particular caller, it is impossible to determine which
implementation is in use without inspecting the code
\footnote{
	Of course, a caller who wanted to use different keywords in their dictionary would
	notice the difference between these implmentations very quickly!
}.

However, it is likely that using named values for a function which does not accept them
will result in an error.
This is because names are themselves values and are unlikely to be the correct kind of
value for the function.
For example, this attempt to use named values resulted in an error because the local
variable \texttt{a} is bound to the value \texttt{\#:a}, which is not a valid input to the
built-in \texttt{+} procedure:

\begin{pygmented}[lang=lisp]
(define (add-values a b)
  (+ a b))

(add-values #:a 5 #:b 7)
; ice-9/boot-9.scm:1685:16: In procedure raise-exception:
; In procedure +: Wrong type argument in position 1: #:a
\end{pygmented}

\subsubsection{Historical Record}
The Scheme standard requires that implementations accept value positionally.
This has been true since the original paper describing lisp
\parencite[][pages 185-186]{original-lisp-paper-pt-1}.
The stated motivation is to distinguish between "functions" and "forms".
In that paper, a function is the abstract idea of a formula that can have concrete numbers
applied to it.
Traditionally, a form is syntactically identical to a formula, but it is meant to be a
stand-in for whatever value the formula will resolve to in context.
The differentiating syntax is taken from a previous work which had a similar concern and
does not explicitly justify the postional notation
\parencite[][pages 3-7]{calculi-of-lambda-conversion}
\footnote{
	The decision to supply arguments positionally appears to derive from prior work
	describing multi-argument functions as a series of functions which take in one of the
	values and call the "next" function which will receieve the "next" value (this is
	similar to the concept of currying used in some programming languages).
	With this model, it is natural to write values positionally as they must be provided in
	the correct order to ensure that each partial function operates correctly.
	It's also worth noting that the examples of functions used in the paper are relatively
	simple - they are all inlinable with the prose.
	In the face of such simplicity, naming values would be verbose without much benefit.
}.

While the semantics of positional arguments have remained the same, the implementation
has changed significantly.
In the oldest version of Guile available through source control, values are given to
a function by collecting them into a list, then dispatching based on the number of
values expected by the function
\parencite[][libguile/eval.c lines 1822-2004]{guile-first-commit}.
With the addition of the VM, values are given by putting them into a frame managed by the
VM (eg, not a C-level frame), avoiding the need to create a list unless the function is
variadic.

\subsection{Typed Value}
This feature is provided by defining methods in the Guile Object-Oriented Programming
System (GOOPS).

\subsubsection{Syntax \& Semantics}
In GOOPS methods are not contained within classes.
Instead, they are simply functions with type specifications.
When an author defines a method, they can choose to specify a type by providing a
two-element list in place of an argument name
\parencite[][Section 8.6 "Methods and Generic Functons"]{guile-manual}.
The first element of the list is the argument name and the second element is the class
name.
For example:

\begin{pygmented}[lang=lisp]
(define-method (add-values (a <number>) (b <number>))
  (+ a b))
\end{pygmented}

This defines a method that will only be executed when it is called with exactly two
arguments which are both numbers.
The author can also define alternative implementations for different types\footnote{
	While overloading is generally out of scope for this paper, the implementation section
	will not make sense without mentioning it.
}:

\begin{pygmented}[lang=lisp]
(define-method (add-values (a <list>) (b <list>))
  (append a b))
\end{pygmented}

\subsubsection{Implementation}

\footnote{
	In most implementation sections I present the bytecode associated with a function
	definition and call and describe the low-level semantics.
	This would not be useful for this section.
	The GOOPS implementation is primarily concerned with function overloading, not with
	type checking.
	Most of the logic is related to checking the arity of the function and searching for
	the function with the "best fit" type signature (that is, it prefers more specific
	child classes over more general parent classes).
	The fact that it will produce an error if the user supplies invalid types appears to be
	incidental to the system, not a primary purpose.
	In order to avoid a drawn-out explanation which would not be relevant to this paper, I
	refer only to the source implementation which is relevant to type-checking.
}

Defining a method defines two separate functions: a generic which is responsible for
method dispatching and the method which contains the body that the author defined
\parencite[][Section 8.6 "Methods and Generic Functons"]{guile-manual}.
A generic is an applicable struct (a structure which can be called like a function) which
stores all methods which share a name
\parencite[][module/oop/goops.scm lines 2045-2245]{guile-source}.
When a caller uses that name with some arguments, the generic object searches all of the
methods associated with that name for a signature with the best fit
\parencite[][module/oop/goops.scm lines 1381-1449]{guile-source}.
This indirectly causes an error if the caller provides a set of values which do not have a
compatible type: the searching process will fail to find a match and print a message like
this one:

\begin{verbatim}
ice-9/boot-9.scm:1685:16: In procedure raise-exception:
No applicable method for #<<generic> add-values (2)> in call (add-values foo bar)
\end{verbatim}

The bytecode for calling a function with typed values is identical to the simple call.

\subsubsection{Historical Record}
GOOPS was first added to the Guile repository in version 1.6
\parencite[][Section 9.1.4 "A Timeline of Selected Guile Releases"]{guile-manual}.
Its implementation was based on STklos, the object system for STK, and it was also
influenced by CLOS, the object system for Common Lisp
\parencite[][module/oop/goops.scm lines 24-25]{guile-source}
\parencite[][Section 8.0 "GOOPS"]{guile-manual}.
The original code
\footnote{
	I am tentatively using 4b5d86e0334f6b8c0b37c55cf47a4cd30e7803e0 as the commit which 
	"adds GOOPS".
	This is somewhat fictitious; 9 commits prior (fdd70ea97c142dc8db1e3f147ac6f5bd6ae157c6)
	changed the major version "due to the merge of GOOPS".
	However, at the time this commit was made no GOOPS files had been added; the version
	number was adjusted in anticipation of the merge, not after the merge was complete.
	I am choosing the other commit as representing the initial add of GOOPS to the
	repository because it is the last consecutive commit which adds or modifies a GOOPS file
	after the version adjustment.
}
looks quite different from the modern version of the file due to changes in the core
language, bug fixes, and refactorings for performance and readability.
However, the underlying idea of storing all alternative implementations of a function in a
structure that uses a type signature to implementation map has remained stable.

\subsection{Variadic Function}
Provided by the "dotted list" syntax and the \texttt{\#:rest} keyword in a
\texttt{lambda*}.

\subsubsection{Syntax \& Semantics}
The \texttt{\#:rest} syntax allows an author to define a variadic variable that will
contain a (possibly empty) list of variadic values.
The variadic variable must be declared after all other declarations.
For example, this definition and call would have the following value
\footnote{
	As the variadic variable contains a list, the below code uses \texttt{apply} to
	destructure the list.
	This is described in detail in the section on caller destructuring.
}:

\begin{pygmented}[lang=lisp]
(define* (add-values a b #:rest variadic-variable)
	(apply + a b variadic-variable)

(add-values 1000 2001 1002 2003)
; 6006
\end{pygmented}

The dotted list syntax is similar, except that the author uses a single period instead of
the \texttt{\#:rest} keyword and this syntax works without \texttt{optargs}:

\begin{pygmented}[lang=lisp]
(define* (add-values a b . variadic-variable)
	(apply + a b variadic-variable)

(add-values 1000 2001 1002 2003)
; 6006
\end{pygmented}

Other than this syntactic difference, the two methods of defining a variadic function are
identical.

\subsubsection{Implementation}
For definitions, the only notable difference from the simple function is the addition of
the \texttt{bind-rest} instruction near the beginning of the body.
This instruction takes values from the stack and constructs a new list which contains all
of them; this list is bound to the variadic variable
\parencite[][libguile/vm-engine.c lines 756-783,
             libguile/vm.c        lines 1041-1052]{guile-source}.

\begin{pygmented}[lang=c]
static SCM
cons_rest (scm_thread *thread, uint32_t base)
{
  SCM rest = SCM_EOL;
  uint32_t n = frame_locals_count (thread) - base;

  while (n--)
    rest = scm_inline_cons (thread, SCM_FRAME_LOCAL (thread->vm.fp, base + n),
                            rest);

  return rest;
}
\end{pygmented}

For calls, the bytecode is identical to the simple call.

\subsubsection{Historical Record}
The dotted list syntax was first described in the Second Revised Report on Scheme, with
the same semantics that are in use today
\parencite[][page 13]{r2rs}.
The \texttt{\#:rest} syntax was originally used by other lisp dialects and implementations
of Scheme; Guile provided it in the original release of the optargs module in order to
ease the transition for programmers used to those languages
\parencite[][commit 0a852b9424f949575afecb19a391023acc63e635 NEWS lines 849-843]{guile-source}

As described in the section on positional values, the oldest versions of Guile manually
collected argument values into a Scheme-level list and dispatched based on the number of
values given.
This made it trivial for the interpreter to supply the variadic values by dropping
elements from the front of the list of all arguments
\parencite[][libguile/eval.c lines 1652-1655]{guile-first-commit}.

When the VM was added in version 2.0, the \texttt{bind-rest} instruction was included.
It performs essentially the same work as the current implementation
\parencite[][libguile/vm-i-system.c lines 718-732]{guile-2.0}.
